\section{Validation}
\label{sec:valid}

\subsection{One suite, two targets}
The automated suite is written with pytest~\cite{pytest}. It drives the comms
interface with no human interaction and runs against two targets from the same
test files. By default it starts the host
simulator and speaks the protocol over its standard input and output; given
\texttt{-{}-port} it opens a serial port through pyserial~\cite{pyserial} and
speaks to the flashed board instead. The fixture that performs that substitution is the only thing that
differs between the two runs, so a case that passes on the host is the same case
on the board.

\subsection{Thirty-three cases pass on the host build}
The following is a run on the host simulator, on a PC, with no board attached:

{\footnotesize
\begin{verbatim}
$ make -C sim && python -m pytest tests -q
.................................    [100%]
33 passed in 0.10s
\end{verbatim}}

Table~\ref{tab:tests} maps each case the brief requires onto the test that
covers it. The suite is a regression check rather than a demonstration: every
run either reproduces the protocol exactly or fails.

The cases fall into three groups. The first confirms the happy path, that
\texttt{CONNECT} is answered and \texttt{NEW} produces an empty board. The
second confirms that a legal move changes the board \emph{exactly} as expected,
comparing the whole nine-character string rather than the single cell, so a move
that corrupted a neighbouring cell would fail. The third is the authority group:
each test reads the board, issues a request that must be refused, and reads the
board again, asserting the two are identical. Seven of the tests are of that
third kind, which is the property Section~\ref{sec:authority} claims.

Beyond the required cases the suite also covers the robustness paths: an
over-length frame and the recovery after it, non-printable bytes, a request
split across two writes and reassembled, a CRLF terminator, and a checksum that
is correct and one that is not.

\subsection{Results from the board}
\ph{Run the same suite against the flashed board with
\texttt{python -m pytest tests -{}-port /dev/ttyACM0 -q} and paste that output
here, labelled as the board run so it is not confused with the host run above.
Note that pyserial must be installed first.}

\ph{Add the two GUI screenshots to \texttt{figures/} and refer to them here:
one mid-game showing a rejected move and its reason in the status line, and one
after a win showing the cells disabled. Say what the reader should notice in
each.}

Nothing in this section claims a measurement that has not been taken: the
figures quoted above are from the host build, and the board run is recorded
separately once it has been performed.

\begin{table}[htbp]
\caption{Required test cases and where they are covered. Test names are shown without their \texttt{test\_} prefix.}
\label{tab:tests}
\centering
\footnotesize
\begin{tabular}{p{4.0cm}p{3.3cm}}
\hline
\textbf{Required by the brief} & \textbf{Test} \\
\hline
CONNECT returns OK & \texttt{connect\_\allowbreak returns\_\allowbreak ok} \\
NEW returns an empty board & \texttt{new\_\allowbreak returns\_\allowbreak an\_\allowbreak empty\_\allowbreak board} \\
A legal move changes the board & \texttt{a\_\allowbreak legal\_\allowbreak move\_\allowbreak changes\_\allowbreak the\_\allowbreak board\_\allowbreak exactly} \\
Occupied cell rejected, board unchanged & \texttt{occupied\_\allowbreak cell\_\allowbreak is\_\allowbreak rejected\_\allowbreak and\_\allowbreak nothing\_\allowbreak changes} \\
Out of range rejected & \texttt{out\_\allowbreak of\_\allowbreak range\_\allowbreak is\_\allowbreak rejected} \\
Move after game over refused & \texttt{moves\_\allowbreak after\_\allowbreak game\_\allowbreak over\_\allowbreak are\_\allowbreak refused} \\
Scripted win and draw & \texttt{a\_\allowbreak win\_\allowbreak is\_\allowbreak detected}, \texttt{a\_\allowbreak draw\_\allowbreak is\_\allowbreak detected} \\
\hline
\end{tabular}
\end{table}
